\documentclass[11pt]{article}
\usepackage{fontspec,amsmath,amssymb,amsthm,hyperref}
\usepackage[margin=1in]{geometry}
\setlength{\emergencystretch}{3em}
\title{The normalized kinetic Gibbs candidate and sharp parameter concentration}
\author{LangevinLean mathematical exposition}
\date{}
\begin{document}\maketitle
\noindent Source module: \texttt{KIPLMC/Target.lean}. Compile this proof document with XeLaTeX or LuaLaTeX. The accompanying HTML pairs the mathematical argument with the exact Lean statements and records the scope of the formalization.


The velocity-first phase target is the product of an explicit Gaussian thermal velocity measure and the already normalized lifted Gibbs position measure. Its parameter marginal is the powered likelihood, and its concentration radius has the paper's sharp dimension coefficient. These results concern a concrete probability measure; the module deliberately does not claim invariance of that measure under a kinetic diffusion.


\section{Construct thermal velocities as a Gaussian pushforward}


Start with the standard Gaussian on the scaled particle Euclidean space and map \(v\mapsto v/\sqrt N\). \texttt{velocityTarget} defines this measure, giving covariance \(I/N\) for positive \(N\). The probability proof uses measurability of the scalar map, evaluates its preimage of the whole space, and obtains mass one. The definition also makes sense at \(N=0\) under Lean's total inverse convention, but the physical target and subsequent normalization theorems require \(N>0\). No unspecified Gaussian velocity law is assumed.



\noindent Lean declarations: \texttt{velocityTarget}, \texttt{velocityTarget\_probability}.

\section{Normalize the product and compute its position marginal}


\texttt{gibbsPhaseTarget} is \(\mathcal N(0,I/N)\otimes\pi_N\), with velocity first. Joint strong convexity and differentiability of the original pair potential imply that the shared lifted Gibbs target \(\pi_N\) is a probability measure. Lean installs that proved instance and infers probability of the product. Mapping the product by its second projection integrates out the first factor, whose total mass is one, so the result is exactly \(\pi_N\). The proof uses \texttt{Measure.map\_snd\_prod} and simplifies the mass scalar. This computes the normalized marginal rather than merely comparing unnormalized density expressions.



\noindent Lean declarations: \texttt{gibbsPhaseTarget}, \texttt{gibbsPhaseTarget\_probability}, \texttt{gibbsPhaseTarget\_position}.

\section{Identify the actual powered likelihood parameter marginal}


Map the phase target to \(\theta\) by composing position projection with the parameter continuous linear map. The pushforward composition theorem identifies this with first taking the position marginal and then the parameter marginal. The preceding result gives the exact lifted Gibbs position law; the shared IPLA theorem identifies its parameter marginal as \(p_\theta(y)^N\) after normalization. Measurability of both maps is provided explicitly. The theorem therefore connects the concrete phase target with the same normalized marginal target used by the concentration proof.



\noindent Lean declarations: \texttt{gibbsPhaseTarget\_parameter}.

\section{Construct the unique marginal maximizer and its sharp second moment}


Joint strong convexity makes the negative log marginal strongly convex after latent integration. The shared IPLA theorem supplies existence and uniqueness of its minimizer, equivalently of the marginal likelihood maximizer. For that \(\theta_*\), the checked sharp concentration theorem gives \(\int\|\theta-\theta_*\|^2d\pi_\Theta\le d_\theta/(Nm)\). The declarations reuse these proved theorems for the actual target, not an abstract measure with a stated concentration field. The last theorem expresses the square-root bound in the existing extended-nonnegative transport API as \(\operatorname{projectedRMS}(\pi_N,P_\theta,\theta_*)\le\sqrt{d_\theta/(Nm)}\).



\noindent Lean declarations: \texttt{existsUnique\_mmle}, \texttt{parameter\_secondMoment\_le}, \texttt{parameter\_concentration}.

\end{document}
